% =============================================================================
% SCRATCH NOTES — Experiments, analyses, and open questions for LLAMIA
% =============================================================================

% =============================================================================
% CRITICAL EXPERIMENTS NEEDED FOR THE NEW NARRATIVE
% =============================================================================

% 1. RL-OPTIMIZED TOOL CALLING BASELINE (THE KEY NEGATIVE RESULT)
%    - Train GRPO on text-based tool calling: same training setup, same GRPO,
%      but agent communicates through text (top-k moves + scores) instead of
%      latent tokens
%    - Must demonstrate THREE things:
%      (a) GRPO tool calling >> vanilla tool calling (shows RL helps)
%      (b) GRPO tool calling << internalization (shows channel is bottleneck)
%      (c) Gap scales with task horizon (multi-step >> single-step gap)
%    - This is THE differentiating experiment. Without it, the paper is a
%      method proposal. With it, it's an investigation with a finding.
%    - Expected results: ~60-70% of internalization on multi-step tasks,
%      ~85-90% on single-step tasks
% \somesh{Yes, we do it, infact we ablate by adding noise to the latent tokens}

% 2. BOTTLENECK SCALING ANALYSIS (Figure~\ref{fig:bottleneck})
%    - Plot: (internalization - tool calling) performance gap vs task horizon
%    - X-axis: task horizon (move prediction=1 step, behavior cloning=few,
%      commentary=many, instruction planning=most)
%    - Y-axis: normalized performance gap
%    - Should show monotonically increasing gap
%    - This is the money figure — directly visualizes the bottleneck compounding

% 3. DIFFICULTY ESTIMATION (EMERGENT CAPABILITY)
%    - Data: Lichess puzzle ratings + human solve rates stratified by Elo
%    - Task: Given a position, predict difficulty for players at Elo X
%    - Why internalization helps: difficulty depends on the SHAPE of the move
%      distribution — how many "reasonable" alternatives exist, how close they
%      are in evaluation. Text gives top-3 moves; latent gives the full policy.
%    - Metric: Spearman correlation with human difficulty ratings
%    - Key control: show text-based system CANNOT do this well because it
%      lacks distributional information

% 4. BEHAVIORAL STYLOMETRY (EMERGENT CAPABILITY)
%    - Data: Lichess player games (players with 100+ games each)
%    - Task: Given N games, identify player's skill level, style
%      (aggressive/positional/tactical), or match games to player
%    - Why internalization helps: style = systematic deviation from optimal
%      policy. Requires comparing human moves against full policy distribution.
%      Text interface loses the "road not taken" (moves considered but not played)
%    - Metrics: Elo prediction MAE, style classification accuracy
%    - Compare against: text-based (limited to top-k engine outputs),
%      Maia-style dedicated models

% 5. INFORMATION-THEORETIC BOTTLENECK MEASUREMENT *** NOW CRITICAL ***
%    *** ALL 4 REVIEWERS DEMANDED THIS. NOT OPTIONAL. ***
%    - Directly measure: MI(agent_repr; text_summary) vs MI(agent_repr; latent_tokens)
%    - Proxy: train probes on text summaries vs latent tokens to predict
%      various board/game properties; gap in probe accuracy = information loss
%    - Properties to probe: piece placement, king safety, pawn structure,
%      tactical motifs, material balance, player to move advantage
%    - This would give a BITS measurement of the bottleneck — very compelling
%      for the "information bottleneck" narrative
%    - R2: "Calling it 'concrete and measurable' without actually measuring
%      it in information-theoretic terms is overclaiming"
%    - R3: "Even a simple probe comparison would substantially strengthen
%      the claim"

% 5b. "MAXIMAL TEXT" BASELINE *** NOW CRITICAL ***
%    *** R2 + R4 both flagged the bandwidth confound as the #1 technical issue ***
%    - Serialize FULL policy distribution as structured text (all ~30 moves
%      with probabilities, value estimates, top features)
%    - Must show internalization still wins even when text gets the SAME info
%    - This isolates channel FORMAT from channel BANDWIDTH
%    - Without this, central claim is underdetermined (R4: "impossible to
%      disentangle 'text is lossy' from 'the particular text summarization
%      chosen is lossy'")
%    - Also test: multi-turn querying where LLM can request additional info
%    - The intro now PROMISES this experiment (P3 + P5 + results bullet 1)

% =============================================================================
% ADDITIONAL ANALYSES (strengthen existing contributions)
% =============================================================================

% 6. CROSS-AGENT TRANSFER
%    - Does the adapter work across Lc0 network configs? (BT4 vs BT3 vs T80)
%    - Even showing 2-3 configs strengthens "general method" claim
%    - Stretch: can we internalize Stockfish NNUE representations?
%      (different architecture, different representation space)

% 7. ROUTING ANALYSIS
%    - When does the model choose which agent config?
%    - Hypothesis: routes to stronger configs for tactical positions,
%      weaker for positional/quiet positions (matching human intuition)
%    - Visualize: routing decisions overlaid on position characteristics
%    - Important for the "RL-trained invocation" contribution

% 8. K (NUM TOKENS) SENSITIVITY
%    - How does performance vary with K? (K=4, 8, 16, 32, 64)
%    - Does more tokens always help? Is there diminishing returns?
%    - At what K does internalization match text-based? (upper bound on text info)

% 9. TOKEN CONTENT ANALYSIS
%    - What do the 32 latent tokens encode?
%    - Linear probes on individual tokens → which encodes policy, value, features?
%    - Does the model learn to "read" specific tokens for specific info?
%    - Connects to interpretability literature (chess as model system)

% =============================================================================
% LLAMIA-BENCH TASK DEFINITIONS
% =============================================================================

% Category 1: DECISION MAKING
%   - Gameplay (Elo measurement against known engines)
%   - Behavior cloning (move prediction across Elo buckets: 1000-2500)

% Category 2: BEHAVIORAL MODELING
%   - Stylometric profiling (identify player from game characteristics)
%   - Difficulty estimation (predict human success rate on positions by Elo)
%   - Interestingness estimation (predict which positions humans find engaging)

% Category 3: COMMUNICATION
%   - Move commentary (explain individual moves, G-eval + human eval)
%   - Game commentary (summarize/analyze full games)

% Category 4: PLANNING
%   - Instruction-conditioned planning (play to achieve specific goals:
%     "trade into a won endgame", "maintain pressure on the kingside")

% Category 5: STATE UNDERSTANDING
%   - Position evaluation (assess who is winning and why)
%   - Board comprehension (answer questions about the current state)

% =============================================================================
% ANTICIPATED REVIEWER OBJECTIONS (from 16-paper review analysis)
% =============================================================================

% Q: "Will this work beyond chess?"
% A: Method requires only: (1) pretrained agent with extractable representations,
%    (2) lightweight adapter, (3) GRPO training on augmented traces.
%    All three are domain-agnostic. Chess enables MEASUREMENT of the bottleneck
%    because agent representations are well-characterized via interp work.
%    The bottleneck itself is domain-agnostic.
%    NOTE: Do NOT claim chess performance predicts general reasoning.
%    DO claim: the insight (text is the bottleneck) and the method (internalization)
%    transfer. Chess provides the controlled setting to PROVE this.

% Q: "This is a chess paper, not an ML paper"
% A: Chess is the testbed, not the contribution. The contributions are:
%    (1) systematic analysis establishing text as the bottleneck (not invocation)
%    (2) a general method for closing the bottleneck
%    (3) a benchmark for studying heterogeneous MAS
%    Cite chess-as-model-system literature: Chessformer, DiffuSearch,
%    Transcendence, Lc0 look-ahead, Creative Puzzles.

% Q: "Why not test on another domain?"
% A: Measuring the bottleneck requires knowing what the agent's representations
%    contain (to quantify what text discards). Chess is the only domain where
%    interpretability work has mapped these representations. Testing on another
%    domain would show the method works but couldn't explain WHY.
%    Future work: robotics (RT-2), reward models, scientific simulation.

% Q: "Is this just multimodal projection (LLaVA)?"
% A: Unlike static projection, our approach uses DYNAMIC re-encoding — fresh
%    tokens at each state transition. The InvokeAgent ablation (64% vs 1-4%)
%    proves this is the critical differentiator. LLaVA-style projection
%    would degrade to "static multimodal projection" (our P3 baseline).

% Q: "The negative result (tool calling fails) — isn't that obvious?"
% A: No. RL-optimized tool calling (GRPO) substantially improves over naive
%    approaches. The non-obvious finding is that this improvement has a ceiling:
%    no amount of invocation strategy optimization compensates for channel loss.
%    This is analogous to discovering that no amount of compression can recover
%    information lost at the source.

% =============================================================================
% TITLE ALTERNATIVES (for discussion)
% =============================================================================

% Current: "Internalizing Heterogeneous Agents in Language Model Reasoning"
%   + "Heterogeneous" signals non-language agents specifically
%   + Method-forward but not chess-forward
%   - Doesn't foreground the investigation/question angle

% Alt 1: "Beyond Text: The Information Bottleneck in LLM-Agent Collaboration"
%   + Question-forward
%   - Vague, could mean many things

% Alt 2: "The Text Bottleneck: Why Tool Calling Fails for Neural Agent Collaboration"
%   + Provocative, foregrounds the negative result
%   - Might alienate tool-calling community

% Alt 3: "From Tools to Teammates: Internalizing Neural Agent Representations
%         in Language Model Reasoning"
%   + Captures the tools-vs-agents distinction
%   - Long

% RECOMMENDATION: Keep current title. It's precise and appropriately scoped.

% =============================================================================
% STRUCTURAL SUGGESTIONS (for methodology/experiments sections)
% =============================================================================

% The experiments section should be restructured to match the new narrative:
%
% Section X.1: LLAMIA-Bench (task definitions, metrics, data)
% Section X.2: Baselines
%   - Language-only (ChessGPT-style fine-tuning)
%   - Vanilla tool calling (text interface, no RL)
%   - RL-optimized tool calling (GRPO, text interface)  ← NEW, critical
%   - Static projection (LLaVA-style, no dynamic re-encoding)
%   - Task-specific specialists (Maia, dedicated engines)
% Section X.3: Main Results (Table spanning all tasks x all methods)
% Section X.4: Bottleneck Analysis
%   - Gap vs task horizon (Figure~\ref{fig:bottleneck})
%   - Information-theoretic measurement (if available)
% Section X.5: Ablations
%   - InvokeAgent removal (existing)
%   - Stage 1 vs Stage 2 (existing)
%   - K sensitivity (new)
%   - Routing analysis (new)
% Section X.6: Emergent Capabilities
%   - Difficulty estimation (new)
%   - Behavioral stylometry (new)

% =============================================================================
% SIMULATED REVIEWER FEEDBACK (4 NeurIPS reviewers, 2026-04-23)
% Scores: R1=5, R2=5, R3=6, R4=5 → borderline reject
% =============================================================================

% WHAT ALL 4 REVIEWERS AGREED ON:
% 1. The research QUESTION is important and well-motivated
% 2. The ablation (64% vs 1-4%) is the strongest result
% 3. Need formal info-theoretic bottleneck measurement
% 4. Need "maximal text" baseline controlling for bandwidth
% 5. Generalization claims need hedging (done in revision)
% 6. Dynamic re-encoding is the genuine novelty, not the adapter

% UNIQUE CONCERNS BY REVIEWER:
% R1 (generalizability): "Three requirements not unique to chess" → ADDRESSED
%     by strengthening interpretability as the key differentiator (revision P4)
% R2 (technical): "Is this bandwidth or format?" → ADDRESSED by promising
%     maximal-text baseline; NEED TO ACTUALLY RUN IT
% R3 (big picture): "Reach > grasp but question matters" → best ally,
%     gave 6/10. Address empty citations + TODOS (done)
% R4 (MAS expert): "CommNet/TarMAC already studied this" → ADDRESSED by
%     distinguishing co-trained MARL from independently-pretrained LLM setting
%     (revision P2). Also: "GRPO text experiment may not exist yet" → MUST RUN

% CITATIONS NEEDED (from R4 + general):
% \cite{sukhbaatar2016learning} — CommNet: Learning Multiagent Communication
% \cite{jiang2018learning} — ATOC/Learning Attentional Communication
% \cite{das2019tarmac} — TarMAC: Targeted Multi-Agent Communication
% \cite{diffusearch2025} — DiffuSearch: implicit search via diffusion
% \cite{searchlesschess2024} — Searchless Chess / Transcendence
% \cite{mcilroyyoung2020maia} — Maia: human-like chess engine
% \cite{karvonen2024chessformer} — Chessformer architecture

% =============================================================================
% PATH TO ACCEPTANCE: PRIORITIZED ACTION ITEMS
% After 2 rounds of simulated review: 6/5/6/5 → need 7/7/7/6
% =============================================================================

% TIER 1 — MUST HAVE (blockers for all reviewers):
%
% (a) MAXIMAL-TEXT BASELINE *** THE SINGLE MOST IMPORTANT EXPERIMENT ***
%     - Serialize full Lc0 policy (all ~30 move probs), value head, top
%       features as structured JSON/text
%     - Train with same GRPO setup, same compute budget as LLAMIA
%     - Show internalization still wins, especially on multi-step tasks
%     - THIS PROVES "format not bandwidth" — currently just a promise
%     - R2: "I cannot raise when the central claim lacks supporting data"
%     - R3: "This alone would likely be sufficient to move me to 7"
%     - R4: "impossible to disentangle without this baseline"
%
% (b) GRPO TEXT TOOL-CALLING BASELINE (distinct row in main table)
%     - Must show: GRPO text >> vanilla text (RL helps) but << internalization
%     - R4 caught that this experiment may not exist yet
%     - The negative result the entire narrative is built on
%
% (c) FIX K=32 vs K=16 INCONSISTENCY
%     - Abstract/intro say K=32, methodology says K=16
%     - R2 and R4 both flagged this — undermines confidence

% TIER 2 — STRONGLY RECOMMENDED (moves R2→7, R4→6):
%
% (d) INFO-THEORETIC PROBES
%     - Linear probes on latent tokens vs text embeddings predicting:
%       piece placement, king safety, pawn structure, tactical motifs
%     - Report accuracy gap = bits-level bottleneck measurement
%     - Grounds "concrete and measurable" claim formally
%
% (e) BOTTLENECK SCALING FIGURE (Figure~\ref{fig:bottleneck})
%     - Gap (internalization - text) vs task horizon
%     - Intro REFERENCES this figure — it must exist
%     - Ordered: move prediction < behavior cloning < commentary < planning

% TIER 3 — NICE TO HAVE (moves R1→7+, R4→7):
%
% (f) Surface cross-domain results in intro (image pref + persuasion)
% (g) Non-chess pilot experiment (Atari/gridworld — even negative is useful)
% (h) Shared-workspace MAS literature (ChatDev, Generative Agents)
% (i) Gradient flow clarification: does GRPO backprop through adapter?
%     If yes → partial co-training, address honestly
%     If no → explain how model learns to use latent tokens
% (j) Token content analysis / what do the K tokens encode
% (k) Dedicated "Bottleneck Analysis" subsection in experiments
%
% ESTIMATED SCORE IMPACT:
% Tier 1 alone:        6/5/6/5 → 6/6/7/6 (borderline accept)
% Tier 1 + Tier 2:     → 7/7/7/6 (solid accept)
% Tier 1 + 2 + 3:      → 7/7/8/7 (strong accept)
